\documentclass[a4paper,12pt,twoside]{article}
\usepackage[utf8]{inputenc}
\usepackage{graphicx}
\usepackage{subfig}
\usepackage[table,xcdraw]{xcolor}
\usepackage{amsmath,amssymb,amsfonts,bm}
\usepackage[bottom=2cm,top=2cm,left=2.0cm,right=2.0cm]{geometry}
\usepackage{float}
\usepackage[portuges]{babel}
\usepackage{indentfirst}
\usepackage[colorlinks=true, allcolors=black]{hyperref}
\usepackage[alf]{abntex2cite}
\usepackage[export]{adjustbox}
\usepackage{afterpage}
\usepackage{multicol}
\usepackage{multirow}
\usepackage{textcomp}
\usepackage{ dsfont }
\usepackage{listings}

\usepackage{fancyhdr}
\pagestyle{fancy}
\fancyhf{}
\fancyhead[RO,LE]{\nouppercase{\emph\leftmark}}
\fancyhead[LO,RE]{\emph \thepage}
\renewcommand{\sectionmark}[1]{\markboth{#1}{}}

\newtheorem{teo}{Teorema}[section]
\newtheorem{lema}[teo]{Lema}
\newtheorem{cor}[teo]{Corolário}
\newtheorem{prop}[teo]{Proposição}
\newtheorem{defi}{Definição}
\newtheorem{exem}{Exemplo}

\usepackage{color}

\definecolor{codegreen}{rgb}{0,0.6,0}
\definecolor{codegray}{rgb}{0.5,0.5,0.5}
\definecolor{codepurple}{rgb}{0.58,0,0.82}
\definecolor{codeyellow}{rgb}{0.67,0.67,0.0}
\definecolor{backcolour}{rgb}{0.95,0.95,0.92}

\lstset{
language=R,   % R code
literate=
{á}{{\'a}}1
{à}{{\`a}}1
{ã}{{\~a}}1
{é}{{\'e}}1
{ê}{{\^e}}1
{í}{{\'i}}1
{ó}{{\'o}}1
{õ}{{\~o}}1
{ú}{{\'u}}1
{ü}{{\"u}}1
{ç}{{\c{c}}}1
}
\lstdefinestyle{mystyle}{
    language=R,
    backgroundcolor=\color{backcolour},   
    commentstyle=\color{codegreen},
    keywordstyle=\color{blue},
    emphstyle=\color{codeyellow},
    numberstyle=\scriptsize\color{black},
    stringstyle=\color{codegreen},
    basicstyle=\rmfamily\scriptsize,
    breakatwhitespace=false,         
    breaklines=true,                 
    captionpos=b,                    
    keepspaces=true,                 
    numbers=left,                    
    numbersep=6pt,                  
    showspaces=false,                
    showstringspaces=false,
    showtabs=false,                  
    tabsize=2
}

\lstset{style=mystyle}

\begin{document}
\newgeometry{bottom=2.0cm,top=1cm,left=2.0cm,right=2.0cm}
\thispagestyle{empty}
	\begin{figure}[htb!]
		\begin{flushright}
			\includegraphics[scale=.3]{UNICAMP_logo.jpg} 
		\end{flushright}
	\end{figure}
	\vspace{-3.5cm}
	\hspace{1.5cm}
	\begin{flushleft}
	\begin{minipage}{15cm}
	\textbf{IMECC-UNICAMP\\
	Atividade 1 - Métodos computacionais}\\
	Professor: Dr. Guilherme Ludwig\\
	Vitor Macedo Rocha RA: 216394
	\end{minipage}
	\end{flushleft}
\noindent\rule{17cm}{0.4pt}

\section{Distribuições condicionais completas e implementação Gibbs}

Considere $\mathbf{p}=[p_{1},...,p_{14}]^{\top}$, $\mathbf{x}=[n_1,...,n_{14},m_1,...,m_{14}]^{\top}$ e $\mathbb{B}=\{0,1,2,...,N\}$, temos que a verosimilhança e as distribuições a priori são dadas por
\begin{align}
&\mathcal{L}(N,\mathbf{p}|\mathbf{x})\propto \frac{N!}{(N-r)!}\prod_{i=1}^{14}p_{i}^{n_i}(1-p_i)^{N-n_i}\mathds{I}(N \in \mathbb{N})\mathds{I}(p_i \in [0,1])\mathds{I}(r \in \mathbb{B})\\
&\pi(N)=\frac{e^{-\lambda}\lambda^{N}}{N!} \mathds{I}(N \in \mathbb{N})\\
&\pi(p_i)=\mathds{I}(p_i \in [0,1]) , i=1,2,...,14.
\end{align}

A seguir vamos derivar a distribuição condicional de $N|\mathbf{p,x}$. Note que $\mathds{I}(N \in \mathbb{N})\mathds{I}(p_i \in [0,1])\mathds{I}(r \in \mathbb{B})=\mathds{I}(N \in {r,r+1,...})\mathds{I}(p_i \in [0,1])$, pois $\mathds{I}(r \in \mathbb{B})=\mathds{I}(N \in \{r,r+1,...\})$.

\begin{align*}
\pi(N|\mathbf{p,x})&\propto \mathcal{L}(N,\mathbf{p}|\mathbf{x}) \pi(N)\\
&= \frac{N!}{(N-r)!}\prod_{i=1}^{14}p_{i}^{n_i}(1-p_i)^{N-n_i}\mathds{I}(N \in \{r,r+1,...\})\mathds{I}(p_i \in [0,1])\frac{e^{-\lambda}\lambda^{N}}{N!} \mathds{I}(N \in \mathbb{N})\\
&=\frac{e^{-\lambda}\lambda^{N}}{(N-r)!}\prod_{i=1}^{14}p_{i}^{n_i}(1-p_i)^{N}(1-p_i)^{-n_i}\mathds{I}(N \in \{r,r+1,...\})\mathds{I}(p_i \in [0,1])\\
&=\frac{e^{-\lambda}\lambda^{N}}{(N-r)!}\prod_{i=1}^{14}\left[p_{i}^{n_i}\right]\prod_{i=1}^{14}\left[(1-p_i)^{N}\right]\prod_{i=1}^{14}\left[(1-p_i)^{-n_i}\right]\mathds{I}(N \in \{r,r+1,...\})\mathds{I}(p_i \in [0,1])\\
& \propto \frac{e^{-\lambda}\lambda^{N}}{(N-r)!}\left[\prod_{i=1}^{14}(1-p_i)\right]^{N}\mathds{I}(N \in \{r,r+1,...\})\mathds{I}(p_i \in [0,1])\\
& = \frac{e^{-\lambda}\left[\lambda\prod_{i=1}^{14}(1-p_i)\right]^{N}}{(N-r)!}\mathds{I}(N \in \{r,r+1,...\})\mathds{I}(p_i \in [0,1])\frac{e^{-\lambda\prod_{i=1}^{14}(1-p_i)}}{e^{-\lambda\prod_{i=1}^{14}(1-p_i)}}\\
& \propto \frac{\exp\{-\lambda\prod_{i=1}^{14}(1-p_i)\}\left[\lambda\prod_{i=1}^{14}(1-p_i)\right]^{N}}{(N-r)!}\mathds{I}(N \in \{r,r+1,...\})\mathds{I}(p_i \in [0,1])\\
& \propto \frac{\exp\{-\lambda\prod_{i=1}^{14}(1-p_i)\}\left[\lambda\prod_{i=1}^{14}(1-p_i)\right]^{N-r}}{(N-r)!}\mathds{I}(N \in \{r,r+1,...\})\mathds{I}(p_i \in [0,1])\\
& = \frac{\exp\{-\lambda\prod_{i=1}^{14}(1-p_i)\}\left[\lambda\prod_{i=1}^{14}(1-p_i)\right]^{N-r}}{(N-r)!}\mathds{I}(N-r \in \mathbb{N})\mathds{I}(p_i \in [0,1])\\
\end{align*}
Portanto
\begin{equation}\label{piN}
\pi(N|\mathbf{p,x})\propto \frac{\exp\{-\lambda\prod_{i=1}^{14}(1-p_i)\}\left[\lambda\prod_{i=1}^{14}(1-p_i)\right]^{N-r}}{(N-r)!}\mathds{I}(N-r \in \mathbb{N})\mathds{I}(p_i \in [0,1])
\end{equation}

Temos que $(N-r)|(\mathbf{p,x})\sim Poisson(\lambda\prod_{i=1}^{14}(1-p_i))$, identificado pelo \textit{kernel} apresentado em \ref{piN}.
\newpage
\restoregeometry

\newpage
\section{Implementação \textit{Hamiltonian} Monte Carlo}

\newpage
\section{Implementação \textit{Hamiltonian} Monte Carlo no STAN}
\end{document}
























